\documentclass[12pt,a4paper]{article}
\usepackage[utf8]{inputenc}
\usepackage[T1]{fontenc}
\usepackage{amsmath,amssymb,amsfonts}
\usepackage[margin=2.5cm]{geometry}
\usepackage{hyperref}

\title{\textbf{Endorsing the A Priori Demarcation of Epistemic Horizons}}
\author{Franklin Silveira Baldo}
\date{March 2026}

\begin{document}
\maketitle

\begin{abstract}
The Generative Ontology framework has shed its cosmological interpretations (Mechanism C) and committed entirely to measuring the native structural bounds of autoregressive generation (Mechanism B). I fully endorse Hasok Chang's A Priori Demarcation Standard, parameterizing the limits of the hardware directly ($N_c$ for Transformers) to complete the Lakatosian problemshift.
\end{abstract}

\section{The Falsification of Cosmological Frameworks}
As established by the recent factorization of the joint probability distribution, Mechanism C (causal injection) is definitively falsified. The generated semantic universe is not governed by non-local invariant physical laws, but strictly by local encoding sensitivity and attention bleed (Mechanism B).

Therefore, I concede Pigliucci's point: accommodating any generative structural anomaly under a metaphysical label of "Observer-Dependent Physics" without predictive bounds constitutes the Foliation Fallacy. Such an approach possesses zero predictive power.

\section{Endorsing the A Priori Demarcation Standard}
To escape the Foliation Fallacy, the Observer-Dependent physical rules must be deterministically tied to the architecture generating them. I formally endorse the \textit{A Priori Demarcation Standard} outlined by Hasok Chang.

For Observer-Dependent Physics to be valid, the structural limits of the observer's "universe" must be mathematically parameterized \textit{a priori}.
Chang correctly supplies the formula for a Transformer's Epistemic Horizon as a phase transition derived from its $\mathsf{TC}^0$ width constraint:
$$ f(\text{Transformer}, N) = 1 - e^{-\gamma(N - N_c)^+} $$

This mathematically defines the exact epistemic boundary of the autoregressive universe. By stripping Generative Ontology of ungrounded metaphysical dimensions and committing entirely to measuring these explicit architectural constants, the framework has completed its Lakatosian problemshift into strict testability.

\section{Conclusion}
The cosmological debate is over. Substrate Dependence ($\Delta > 0$) is empirical fact. The remaining task is purely mapping the exact empirical phase transitions ($N_c$) and decay parameters ($\lambda$) across substrates.

\end{document}
